% ---------------------------------------------------------------------------------------------
% DRAFT appendix: finite-difference step-size convergence test (refactor plan Step C3).
%
% \input{} at the end of whitepaper.tex.
%
% Figure dependency: c3_step_sweep.png, produced by
%     ./fishertest --sweep > c3_sweep.csv
%     python tests/c3_step_sweep.py c3_sweep.csv -o c3_step_sweep.png
% Copy it into Whitepaper/ before compiling.
% ---------------------------------------------------------------------------------------------

\appendix

\section{Finite-difference step sizes and the convergence of the Fisher forecast}
\label{app:stepsize}

Every derivative in Section~\ref{sec:fisher} is computed numerically, by perturbing a parameter
and re-evaluating the model. The forecast therefore inherits whatever error that approximation
carries, and the size of the perturbation is a free choice that nothing in the physics fixes.
Because the precision gains this paper reports are in places at the few-per-cent level, a
differencing error at that same scale would be indistinguishable from the signal. This appendix
establishes that it is not.

\subsection{Why a step size has a correct value at all}
\label{app:stepsize_theory}

A central-difference estimate of a derivative,
\[
\frac{\partial m}{\partial\theta} \simeq \frac{m(\theta+h) - m(\theta-h)}{2h},
\]
carries two errors that pull in opposite directions. Expanding in a Taylor series, the neglected
terms give a \emph{truncation} error of order $h^2 m'''$, which grows as the step grows: the
secant through two widely separated points stops matching the tangent. Working against it,
evaluating $m$ in finite precision gives a \emph{round-off} error of order $\epsilon |m| / h$,
which grows as the step shrinks: subtracting two nearly equal numbers cancels their leading
significant digits, and dividing by a small $h$ amplifies what remains. Their sum is U-shaped in
$h$, with a broad flat minimum -- a plateau -- between the two regimes. A trustworthy step lies in
the middle of that plateau, and the plateau's existence, not merely the value chosen from it, is
what has to be demonstrated.

For a Fisher forecast the consequence is sharper than a general loss of accuracy, because the
error does not distribute itself evenly across the matrix. This is developed in
Section~\ref{app:stepsize_eigen}.

\subsection{The measurement}
\label{app:stepsize_measure}

The step for each parameter was swept over twelve decades while all other steps were held at their
default, and the recovered $\sigma$ recorded. The sweep uses a synthetic-event fixture rather than
the full Monte Carlo: seven events spanning $t_E$ from $5$ to $900$~d, each in two versions --
peaking inside a Roman high-cadence season and peaking in a Roman gap -- evaluated separately for
the joint, Rubin-only and Roman-only partitions of Section~\ref{subsec:partition}. Using a fixture is
what makes the test affordable (seconds rather than the hours needed to accumulate a comparable
number of Fisher-scored events from the Monte Carlo), and it isolates the numerical question from
the population statistics.

Two of the nine photometric parameters serve as controls. The model magnitude depends on the
baseline magnitudes $m_{b,0}$ and $m_{b,1}$ linearly and with unit slope, so their finite
difference is exact at any step and their curves must be perfectly flat. Any structure there would
indicate a fault in the sweep machinery rather than in the physics. Both are flat to machine
precision across all twelve decades.

Figure~\ref{fig:stepsweep} shows the result. Both walls of the expected U are visible, bracketing a
plateau roughly seven decades wide, across which $\sigma$ is stable to better than $0.1\%$ for
every parameter and every event.

\begin{figure}[htbp]
	\centering
	\includegraphics[width=\linewidth]{c3_step_sweep}
	\caption{Recovered parameter uncertainty $\sigma$ as a function of the finite-difference step,
	for each of the nine photometric parameters (panels), seven test events (colours) and three
	survey partitions (line styles). The step is plotted relative to the adopted value, so the
	adopted step is the vertical line at $1$ and the shaded band marks the plateau. Curves are
	normalised to the adopted step, so a converged parameter is a flat line at unity. The collapse
	at small steps is round-off; the collapse at large steps is truncation. The two baseline-
	magnitude panels ($m_{b,0}$, $m_{b,1}$) are controls: the model is exactly linear in these, so
	their flatness across the full range validates the sweep itself. Note that both failure modes
	drive $\sigma$ \emph{downward} -- they manufacture apparent information rather than destroying
	it, which is why neither announces itself as an obviously broken result.}
	\label{fig:stepsweep}
\end{figure}

\subsection{The steps adopted here}
\label{app:stepsize_adopted}

The simulation inherited its step sizes from a legacy codebase written for a different
science case, where they were expressed as fixed fractions of each parameter: $0.151$ for $u_0$
(against a typical $u_0 \sim 0.3$), $0.255\,t_E$ for both $t_E$ and $t_0$, $0.251\,\pi_E$ for
$\pi_E$, and $3^\circ$ for the trajectory angle $\xi$. Figure~\ref{fig:stepsweep} places all of
these far up the truncation branch: they are approximately four orders of magnitude larger than
the plateau. Every step is therefore scaled by a single factor of $10^{-4}$, which places all nine
in the middle of the plateau simultaneously. A factor of $10^{-5}$ reproduces every $\sigma$ to
better than $0.3\%$, so the result is not sensitive to where in the plateau the step is taken.

Table~\ref{tab:stepbias} gives the cost of the original choice. The entries are the factor by
which the legacy step \emph{understated} the recovered uncertainty, which is the quantity that
matters: in every case the error is in the optimistic direction.

\begin{table}[htbp]
	\centering
	\caption{Factor by which the legacy finite-difference steps understated the forecast
	uncertainty, $\sigma_{\mathrm{converged}} / \sigma_{\mathrm{legacy}}$, over all test events and
	survey partitions. A value of $1.00$ means the legacy step was already adequate for that
	parameter. The two baseline magnitudes are the linear controls.}
	\label{tab:stepbias}
	\begin{tabular}{lrr}
		\toprule
		Parameter & Median & Worst case \\
		\midrule
		$u_0$        & 21.7 & 7643 \\
		$t_0$        &  6.0 & 10388 \\
		$t_E$        &  1.45 &  26.1 \\
		$\pi_E$      &  1.00 & 219.8 \\
		$\xi$        &  1.00 &  45.6 \\
		$f_{b,0}$, $f_{b,1}$ & 1.00 & 1.01 \\
		$m_{b,0}$, $m_{b,1}$ (control) & 1.00 & 1.00 \\
		\bottomrule
	\end{tabular}
\end{table}

The median and worst-case columns differ by orders of magnitude because the damage is
concentrated on the events whose Fisher matrices are closest to degenerate -- short events, and
events with sparse coverage. For $\pi_E$ and $\xi$ the median is unity: on a typical
well-sampled event the legacy step was harmless, and the failure appears only where the matrix is
already nearly singular. This is precisely the population that the joint-fit science case is about.

\subsection{Where the error goes: the eigenstructure of the Fisher matrix}
\label{app:stepsize_eigen}

That the errors concentrate on near-degenerate events is not a coincidence, and it explains why
the effect is larger than the per-parameter percentages suggest.

The Fisher matrix is assembled in physical units, so its entries carry units of
$1/(\theta_j\theta_k)$ and span many orders of magnitude before any physics enters --
$t_E$ is measured in days ($\sim 30$), $u_0$ is dimensionless ($\sim 0.3$), $\xi$ is in radians.
Normalising by $\tilde{F} = D F D$ with $D = \mathrm{diag}(1/\sqrt{F_{ii}})$ removes this: the
diagonal becomes unity, every off-diagonal entry becomes a correlation coefficient, and the
inverse is recovered exactly as $F^{-1} = D\tilde{F}^{-1}D$. Whatever ill-conditioning survives in
$\tilde{F}$ is a statement about the data, not about the choice of units.

Diagonalising $\tilde{F}$ rotates onto parameter \emph{combinations} that are measured
independently, with eigenvalue $\lambda_k$ the information carried about combination $k$ and
$\lambda_{\max}/\lambda_{\min}$ the condition number. A large condition number is usually not a
numerical inconvenience but a physical fact: some combination of parameters genuinely cannot be
recovered from that light curve, and the eigenvector belonging to $\lambda_{\min}$ names it.
For a short event observed only sparsely, that direction is the familiar $u_0$--$\pi_E$
degeneracy; for an event with poorly sampled wings it is the $t_E$--$f_b$ degeneracy between a
longer, more blended event and a shorter, less blended one.

Comparing the normalised spectrum of a $5$-day event at the legacy and converged steps makes the
mechanism explicit:

\begin{center}
\begin{tabular}{lcccc}
	\toprule
	& $\lambda_1 \ldots \lambda_6$ & $\lambda_7$ & $\lambda_8$ & $\lambda_9$ \\
	\midrule
	Legacy steps    & $3.30, 1.83, 1.30, 1.19, 0.712, 0.532$ & $7.5\times10^{-2}$ & $6.1\times10^{-2}$ & $5.3\times10^{-3}$ \\
	Converged steps & $3.42, 1.89, 1.30, 1.16, 0.711, 0.513$ & $7.8\times10^{-3}$ & $8.7\times10^{-7}$ & $2.3\times10^{-10}$ \\
	\bottomrule
\end{tabular}
\end{center}

\noindent The six best-constrained directions are essentially unchanged. The entire discrepancy
sits in the three flattest ones, where the smallest eigenvalue is inflated by a factor of
$\sim\!4\times10^{7}$ and the condition number consequently reads $6.2\times10^{2}$ instead of its
true $1.5\times10^{10}$. Since $\tilde{F}$ has unit diagonal its trace is fixed at the dimension,
so no information was created: it was redistributed out of the well-constrained directions into
the degenerate ones. Physically, truncation error was papering over real parameter degeneracies --
the legacy configuration reported a sub-per-cent parallax measurement for an event far too short
to exhibit annual parallax at all.

This is the reason the convergence test is reported here rather than being treated as an internal
numerical detail. An unconverged derivative in a Fisher forecast does not present as scatter; it
presents as a confident, well-conditioned, and entirely spurious measurement of exactly those
parameter combinations the data cannot constrain.

\subsection{Effect on the results of this paper}
\label{app:stepsize_impact}

Two consequences, of very different weight.

First, absolute parameter uncertainties are affected, and every $\sigma$ quoted here is computed
with the converged steps. Forecasts made with the legacy configuration were optimistic, severely so
for short-$t_E$ events.

Second -- and this is what matters for the paper's claims -- the \emph{relative} results are
robust. Because the joint, Rubin-only and Roman-only matrices are accumulated from the same
simulated event with the same random draws, the comparison between them is paired, and a common
multiplicative error in the derivatives largely cancels in the per-event ratio. Table
\ref{tab:stepratio} shows the headline quantity, the ratio of the joint-fit uncertainty on $t_E$ to
the better of the two single-survey fits, computed either side of the correction.

\begin{table}[htbp]
	\centering
	\caption{Ratio $\sigma_{t_E}^{\mathrm{joint}} / \min(\sigma_{t_E}^{\mathrm{Rubin}},
	\sigma_{t_E}^{\mathrm{Roman}})$ for each test event, computed with the legacy and the converged
	finite-difference steps. Values below unity indicate the joint fit outperforms either survey
	alone. Events labelled ``gap'' peak between Roman observing seasons.}
	\label{tab:stepratio}
	\begin{tabular}{lrr}
		\toprule
		Event & Legacy steps & Converged steps \\
		\midrule
		$t_E = 5$~d, in season    & 0.9995 & 0.9981 \\
		$t_E = 5$~d, in gap       & 1.0000 & 1.0000 \\
		$t_E = 25$~d, in season   & 0.9980 & 0.9902 \\
		$t_E = 25$~d, in gap      & 1.0000 & 1.0000 \\
		$t_E = 100$~d, in season  & 0.9935 & 0.9717 \\
		$t_E = 100$~d, in gap     & 0.9788 & 0.9827 \\
		$t_E = 900$~d             & 0.7588 & 0.7951 \\
		\bottomrule
	\end{tabular}
\end{table}

The ordering is preserved and the magnitudes shift by at most a few per cent, with the joint-fit
advantage slightly larger after correction in most cases. The two events where it is smaller are
instructive: in both, the correction falls more heavily on the survey playing the supporting role
than on the one that dominates the fit. For the $100$-day gap event, Roman is the weaker partner --
it barely covers the peak -- and its $\sigma_{t_E}$ inflates by a factor $2.14$ against Rubin's
$1.28$. For the $900$-day event, Rubin's contribution is specifically the parallax channel, its
continuous coverage sampling the annual signal across Roman's seasonal gaps, and its
$\sigma_{\pi_E}$ inflates by $2.18$ against Roman's $1.10$. This mechanism is inferred from seven
constructed events and should be regarded as a plausible reading rather than an established
property.

\subsection{A note on the differencing scheme}
\label{app:stepsize_stencil}

The sweep also settled a question left open in Section~\ref{sec:fisher}: why $t_E$ and $\pi_E$ were
differenced with both evaluations on the same side of the parameter value, rather than
symmetrically. The pattern applies to exactly the parameters constrained to remain positive, which
suggests the intent was to avoid evaluating the model at an unphysical value. That is not a real
constraint at these step sizes, and the cost is significant: averaging two one-sided differences
gives $m'(\theta) + \tfrac{3}{8}m''(\theta)h + \mathcal{O}(h^2)$, a first-order and biased estimate,
against the second-order symmetric form. Verified on a smooth test function, the one-sided scheme
carries $\approx\!22\times$ the error at equal step and improves by only one decade per decade of
step reduction, against two. Both parameters now use the symmetric form. At the converged step the
numerical difference is immaterial, since the bias term scales with $h$; the change removes a
latent sensitivity rather than altering a result.

The corresponding astrometric derivatives ($\theta_E$ and $\pi_E$ in the astrometric matrix) retain
the one-sided scheme, and their steps have not been swept. This is deliberate: the astrometric
forecast also depends on the unresolved question of Roman's F146 astrometric error model noted in
Section~\ref{sec:fisher_caveats}, so its absolute uncertainties are not yet quantitative regardless
of the differencing scheme, and correcting one of two unverified inputs in isolation would not be
checkable. Both should be addressed together.
